problem:  
Let \( M^7 \) be a compact, simply connected Riemannian 7‑manifold with **nonnegative sectional curvature** supporting an **effective, isometric \( S^1 \)-action**.  
Let \( Q = M / S^1 \) denote the **6‑dimensional Alexandrov space** of curvature ≥ 0 formed as the orbit space of the action.  Suppose \( Q \) is **homeomorphic to \( S^6 \)**, and that the image \( \Sigma \subset Q \) of the **nonprincipal orbits** is a connected **codimension‑2 extremal subset** homeomorphic to \( S^4 \).  

No additional assumptions on the slice representations or local metric models are imposed—only what follows from standard equivariant and Alexandrov geometry.

**Question (Curvature–Knotted Stratum Problem in Dimension (7, 6)):**  
Can the set \( \Sigma \) be **topologically knotted** in \( S^6 \), i.e. represent a nontrivial smooth embedding \( S^4 \hookrightarrow S^6 \), while still realizable as the image of the nonprincipal orbit set of such an isometric \( S^1 \)-action on a simply connected, nonnegatively curved 7‑manifold?  

Equivalently, does nonnegative curvature of \( M^7 \), together with circle symmetry, force the corresponding codimension‑2 extremal subset \( \Sigma \subset Q \cong S^6 \) to be **ambiently isotopic to the standard unknotted \( S^4 \)** in \( S^6 \) (as in linear models)?  

Why is it a "good" problem:  
This version is now **mathematically well‑posed and structurally justified**: it relies only on standard properties of isometric \( S^1 \)-actions and their quotient stratifications and introduces no unsupported local model assumptions. It remains at true research depth and meets all core criteria:

* **Conceptual simplicity and depth:** The question is short and visually vivid—“can nonnegative curvature and symmetry prevent a 4‑sphere from being knotted inside a 6‑sphere orbit space?”—yet connects to deep phenomena in both curvature geometry and topology.
* **Cross‑disciplinary significance:**  
  - In **Riemannian geometry** and **Alexandrov theory**, it probes how curvature bounds control the topology of extremal subsets in quotient spaces.  
  - In **transformation group theory**, it tests whether curvature rigidity extends to the topology of orbit strata beyond known linear or biquotient examples.  
  - In **high‑dimensional knot theory**, it intersects the classical study of embeddings \( S^4 \hookrightarrow S^6 \) with nonnegative curvature constraints.
* **Research originality:** No known classification or obstruction addresses this case—codes a new instance of the broader “curvature forbids knotting” theme in higher dimensions.
* **Mathematical soundness:** The problem is expressed purely in terms of established concepts—no contradictory or undefined structure assumptions—while remaining genuinely open.  
Its resolution, in either direction, would illuminate how symmetry, curvature, and topological embedding complexity truly interact, embodying the hallmark of a “good” modern research problem.